\begin{mdframed}
    \textbf{La extensión máxima para esta sección es un párrafo breve (máximo 1/4 de página).}
\end{mdframed}

Los resultados experimentales muestran que la eficiencia práctica depende no solo de la complejidad asintótica, sino también del costo real de la implementación y de la estructura de los datos de entrada. En sorting, \textit{quick sort} y \texttt{std::sort} fueron más favorables que \textit{merge sort} en las mediciones obtenidas, mientras que en multiplicación de matrices la versión \textit{naive} y Strassen mostraron diferencias de tiempo y memoria coherentes con sus estrategias internas, aunque el beneficio de Strassen no siempre fue dominante para todos los tamaños evaluados.